% ============================================================
%  Capítulo -- Planificación y gestión
% ============================================================

\chapter{Planificación y gestión}
\label{sec:planificacion_gestion}

\section{Metodología de trabajo}

\subsection{Adaptación de Scrum}

Para la organización del trabajo se tomó Scrum como referencia metodológica, pero
se aplicó de forma adaptada a las condiciones reales del proyecto. OmniLitter no
se desarrolló dentro de un equipo Scrum completo, sino como un trabajo académico
realizado por dos personas, en coordinación con un cliente externo y bajo la
supervisión del tutor. Por tanto, la metodología se utilizó como marco de trabajo
iterativo e incremental, no como una implantación estricta de todos sus roles,
eventos y artefactos.

La principal adaptación afectó a los roles. El cliente asumió una función próxima
a la de \textit{Product Owner}, ya que aportó la visión de dominio, explicó las
necesidades científicas del sistema, validó los prototipos y ayudó a priorizar las
funcionalidades con mayor valor para el uso real de la plataforma. Esta figura fue
especialmente importante en aspectos como la captura de datos en campo, la
organización por sesiones de muestreo, la taxonomía de residuos, el funcionamiento
sin conexión y la utilidad posterior de los datos exportados. Sin embargo, al
tratarse de un proyecto académico, la responsabilidad final de planificación,
desarrollo y documentación permaneció en el equipo de trabajo.

No existió una figura formal de \textit{Scrum Master}. Las tareas habituales de
seguimiento, eliminación de bloqueos y coordinación se repartieron entre los dos
integrantes del equipo. Esta decisión resultó coherente con el tamaño del proyecto:
introducir un rol específico de facilitación habría sido poco realista para un
equipo tan reducido. En su lugar, se mantuvo una coordinación directa y continua,
centrada en revisar avances, detectar dependencias entre backend, portal web y
aplicación móvil, y ajustar el reparto de tareas según la carga de cada momento.

La planificación se estructuró en sprints, entendidos como bloques de trabajo con
objetivos funcionales concretos. No obstante, su duración no se fijó de manera
completamente rígida. El calendario académico, la disponibilidad de los miembros
del equipo, las reuniones con cliente y tutor, y la complejidad variable de algunas
tareas obligaron a adaptar la duración efectiva de cada iteración. Por este
motivo, los sprints se utilizaron principalmente para ordenar el avance
incremental del producto: primero la base funcional del sistema, después la captura
de campo y el análisis inicial, posteriormente la sincronización, mapas y
exportación, y finalmente el cierre funcional y refinamiento.

También se adaptaron los eventos de Scrum. El \textit{Daily Scrum} no se realizó
como una reunión diaria formal, ya que el equipo estaba formado únicamente por dos
personas. En la práctica, se sustituyó por una coordinación constante mediante
comunicación directa, revisión frecuente del estado de las tareas y resolución
rápida de dudas técnicas. Este formato permitió mantener visibilidad sobre el
progreso sin introducir una ceremonia artificial que no aportaba valor adicional al
tamaño del equipo.

Las revisiones de sprint se realizaron cuando fue posible con la participación del
cliente y del tutor. Estas sesiones permitieron contrastar incrementos funcionales,
validar decisiones de diseño y ajustar el backlog a partir de la retroalimentación
recibida. Cuando no fue viable realizar una revisión formal al cierre exacto de una
iteración, la validación se produjo mediante reuniones de seguimiento, revisión de
prototipos, comprobación de funcionalidades implementadas o discusión de cambios en
los requisitos. De esta forma, se conservó el principio de inspección y adaptación
propio de Scrum, aunque con una ejecución flexible.

En conjunto, la adaptación de Scrum permitió mantener una forma de trabajo
ordenada, incremental y revisable, adecuada tanto al contexto académico como a la
existencia de un cliente real. La metodología ayudó a dividir el desarrollo en
entregas parciales, priorizar funcionalidades según su valor, incorporar cambios de
requisitos sin romper la planificación general y conservar trazabilidad entre
necesidades, backlog, implementación y validación.

\section{Planificación Inicial}

La planificación inicial de OmniLitter tuvo como objetivo transformar los
requisitos obtenidos con el cliente en un conjunto de bloques funcionales
ordenados, asumibles por un equipo de dos personas y coherentes con el calendario
académico. Desde el inicio se identificó que el proyecto no consistía en una única
aplicación aislada, sino en una plataforma formada por tres piezas principales:
backend, portal web y aplicación móvil. Por ello, la planificación se centró en
reducir riesgos de integración y en construir primero las capacidades sobre las
que dependían las funcionalidades posteriores.

La planificación no se entendió como un documento cerrado, sino como una línea de
trabajo inicial revisable. Las reuniones con el cliente, la validación de
prototipos, las restricciones técnicas del entorno móvil y la evolución del alcance
hicieron necesario ajustar prioridades durante el desarrollo. Aun así, se mantuvo
una estructura común: identificar los bloques principales, ordenar sus
dependencias, repartir el trabajo por sprints y revisar los riesgos más relevantes
en cada etapa.

\subsection{Identificación de bloques funcionales}

El primer paso consistió en agrupar los requisitos en bloques funcionales. Esta
agrupación permitió evitar una planificación basada únicamente en tareas pequeñas
y facilitó razonar sobre qué partes del sistema aportaban valor de forma conjunta.
Los bloques identificados fueron los siguientes:

\begin{itemize}[leftmargin=1.4em]
\item \textbf{Autenticación, usuarios y permisos.} Incluye registro, inicio de
sesión, cierre de sesión, roles, grupos de trabajo y control de acceso. Este bloque
era necesario para poder asociar datos a usuarios reales y separar la información
entre grupos.

\item \textbf{Gestión de campañas y sesiones de muestreo.} Define el contexto
organizativo y científico de la captura de datos. Las campañas agrupan el trabajo
de campo y las sesiones representan actividades concretas de recogida.

\item \textbf{Captura de observaciones en aplicación móvil.} Comprende el registro
de residuos, selección de taxonomía, cantidades, notas, coordenadas y fotografías.
Es el bloque central desde la perspectiva del usuario de campo.

\item \textbf{Persistencia local y sincronización.} Permite trabajar sin conexión,
almacenar datos pendientes y sincronizarlos posteriormente con el backend. Se
identificó desde el inicio como una de las áreas de mayor complejidad técnica.

\item \textbf{Consulta, mapa, analítica y exportación.} Agrupa las funcionalidades
del portal web orientadas a revisar, interpretar y reutilizar los datos recogidos:
listados, filtros, visualización geográfica, métricas y exportaciones.

\item \textbf{Administración, configuración y despliegue.} Incluye la gestión
operativa de usuarios, grupos, catálogo científico, preferencias, documentación
técnica y preparación del sistema para su ejecución en entornos reales.
\end{itemize}

Esta división ayudó a mantener la trazabilidad entre requisitos, backlog,
implementación y pruebas. Además, permitió distinguir entre funcionalidades
fundacionales, como autenticación y permisos, y funcionalidades de valor final,
como sincronización, mapas o exportación científica.

\subsection{Dependencias entre módulos}

La planificación inicial tuvo en cuenta que muchos módulos no podían desarrollarse
de forma completamente independiente. El backend actuaba como capa común de datos y
reglas de negocio, mientras que el portal web y la aplicación móvil consumían sus
servicios desde perspectivas distintas. Por este motivo, se priorizó construir
primero los contratos, modelos y permisos compartidos antes de abordar flujos más
avanzados.

\begin{table}[H]
\centering
\caption{Dependencias principales consideradas en la planificación inicial}
\label{tab:dependencias_planificacion_inicial}
\small
\renewcommand{\arraystretch}{1.18}
\begin{tabularx}{\textwidth}{p{3.2cm}X X}
\toprule
\textbf{Módulo o bloque} & \textbf{Dependencias principales} & \textbf{Justificación de planificación} \\
\midrule
Autenticación y permisos & Ninguna dependencia funcional previa relevante. & Debía implementarse al inicio para proteger rutas, asociar datos a usuarios y aplicar visibilidad por grupos. \\
Grupos y membresías & Autenticación. & Los grupos determinan el alcance de campañas, sesiones, observaciones, métricas y administración. \\
Campañas y sesiones & Autenticación, grupos y permisos. & Cada sesión necesita un contexto de grupo, campaña, usuario observador y reglas de acceso. \\
Observaciones y fotografías & Sesiones, taxonomía, permisos y almacenamiento. & Las observaciones solo tienen sentido dentro de una sesión y deben conservar trazabilidad científica. \\
Persistencia offline & Modelos de campañas, sesiones, observaciones y usuario autenticado. & La copia local debía reflejar el modelo del backend para poder sincronizar sin pérdida de información. \\
Sincronización & Persistencia offline, API de sesiones y observaciones, estados locales y remotos. & Era necesario disponer antes de captura local y contratos estables para resolver el envío de cambios pendientes. \\
Portal de consulta y análisis & Observaciones sincronizadas, permisos y filtros por grupo. & El portal depende de datos ya consolidados para listados, métricas, mapa y exportaciones. \\
Despliegue y operación & Backend, web, configuración de entorno y build móvil. & La puesta en marcha requiere que los módulos principales estén integrados y documentados. \\
\bottomrule
\end{tabularx}
\end{table}

Estas dependencias condicionaron el orden de los sprints. En particular, no era
conveniente abordar la sincronización antes de disponer de captura local, ni
desarrollar el análisis web antes de contar con observaciones suficientemente
estructuradas. Del mismo modo, los permisos por grupo debían estar presentes desde
las primeras iteraciones, ya que afectaban transversalmente a casi todos los
módulos.

\subsection{División inicial en sprints}

La planificación se organizó en cuatro sprints principales. Cada sprint se definió
como un incremento con un objetivo funcional reconocible, evitando que las
iteraciones fueran simples agrupaciones técnicas sin valor observable. La duración
real de cada sprint se ajustó a la disponibilidad del equipo y al calendario
académico, pero la secuencia prevista se mantuvo como guía general del desarrollo.

\begin{table}[H]
\centering
\caption{División inicial del trabajo por sprints}
\label{tab:division_inicial_sprints}
\small
\renewcommand{\arraystretch}{1.18}
\begin{tabularx}{\textwidth}{p{2.1cm}p{4.0cm}X}
\toprule
\textbf{Sprint} & \textbf{Objetivo principal} & \textbf{Bloques funcionales previstos} \\
\midrule
Sprint 1 & Base funcional del sistema. & Autenticación, usuarios, grupos, permisos, campañas, sesiones iniciales, taxonomía y primeras observaciones. \\
Sprint 2 & Captura de campo y análisis inicial. & Mejora de sesiones y observaciones, GPS, fotografías, validaciones, persistencia local inicial y consulta web de datos capturados. \\
Sprint 3 & Sincronización, mapas y exportación científica. & Cola de sincronización, publicación de datos offline, visualización geográfica, métricas, exportación CSV/JSON/GeoJSON y preferencias de uso. \\
Sprint 4 & Cierre funcional y refinamiento. & Administración, ajustes de campañas y sesiones, catálogo científico, correcciones finales, pruebas regresivas, documentación y preparación de despliegue. \\
\bottomrule
\end{tabularx}
\end{table}

Esta división permitió construir el sistema de forma incremental. El primer sprint
dejaba una arquitectura mínima conectada; el segundo acercaba la aplicación móvil
al trabajo real de campo; el tercero cerraba el ciclo entre captura offline,
sincronización y explotación de datos; y el cuarto estabilizaba la plataforma para
su entrega y validación final.

\subsection{Riesgos identificados}

Durante la planificación inicial se identificaron varios riesgos que podían afectar
al alcance, al calendario o a la calidad del producto. Estos riesgos se revisaron
durante el desarrollo y explican algunas decisiones de priorización, como construir
primero una base técnica estable, validar pronto el flujo móvil o reservar margen
para integración y despliegue.

\begin{table}[H]
\centering
\caption{Riesgos principales identificados en la planificación inicial}
\label{tab:riesgos_planificacion_inicial}
\small
\renewcommand{\arraystretch}{1.18}
\begin{tabularx}{\textwidth}{p{3.7cm}p{1.8cm}X}
\toprule
\textbf{Riesgo} & \textbf{Impacto} & \textbf{Medidas de mitigación} \\
\midrule
Cambios de requisitos del cliente & Alto & Mantener un backlog revisable, validar prototipos e incrementos funcionales y documentar la evolución de requisitos. \\
Complejidad de sincronización offline & Alto & Desarrollar primero persistencia local, separar estados de sincronización, probar recorridos manuales completos y limitar el alcance a los flujos principales. \\
Coordinación entre portal web y aplicación móvil & Alto & Centralizar reglas en el backend, compartir contratos y esquemas cuando fuese posible, y validar los flujos de extremo a extremo. \\
Curva de aprendizaje de Expo, React Native y emuladores & Medio & Reservar tiempo para configuración, pruebas en emulador y dispositivo real, y priorizar funcionalidades móviles críticas antes del cierre. \\
Problemas de despliegue e infraestructura & Medio & Documentar entornos, utilizar configuración reproducible, preparar despliegue progresivo de backend, web y build móvil, y dejar constancia de limitaciones operativas. \\
Sobrecarga por equipo reducido & Medio & Repartir tareas por módulos, mantener coordinación constante y priorizar funcionalidades imprescindibles frente a mejoras secundarias. \\
\bottomrule
\end{tabularx}
\end{table}

El riesgo más relevante desde el punto de vista técnico fue la sincronización
\textit{offline-first}, porque implicaba coordinar almacenamiento local, estados
pendientes, API remota, reintentos y validación de datos. Desde el punto de vista
de gestión, el riesgo principal fue la evolución de requisitos, propia de un
proyecto con cliente real y dominio científico específico. La planificación trató
ambos riesgos mediante iteraciones cortas, revisiones frecuentes y una separación
clara entre el núcleo funcional imprescindible y las ampliaciones deseables.

\section{Gestión y herramientas de seguimiento}

\subsection{Gestión de la comunicación}

La comunicación del proyecto se organizó diferenciando entre la coordinación
externa, con cliente y tutor, y la coordinación interna del equipo. Para las
reuniones formales y las cuestiones rápidas con el cliente y el tutor se utilizaron
principalmente Microsoft Teams y Outlook. Teams permitió mantener reuniones de
seguimiento, revisar dudas funcionales y compartir aclaraciones sobre requisitos,
mientras que Outlook se empleó para comunicaciones más formales, convocatoria de
reuniones y envío de información relevante.

Entre los dos integrantes del equipo se emplearon canales informales, principalmente
WhatsApp, para coordinar avances diarios, resolver dudas rápidas y avisar de
bloqueos. Esta comunicación directa resultó adecuada para un equipo pequeño, ya
que redujo el coste de coordinación y permitió reaccionar con rapidez ante cambios
de planificación o problemas técnicos.

\begin{figure}[H]
\centering
\fbox{%
\parbox{0.86\textwidth}{%
\centering
\vspace{1.2cm}
Captura pendiente: canales de comunicación utilizados en Microsoft Teams, Outlook
o ejemplo de reunión de seguimiento.
\vspace{1.2cm}
}}
\caption{Herramientas utilizadas para la comunicación con cliente y tutor.}
\label{fig:comunicacion_teams_outlook}
\end{figure}

\subsection{Gestión del backlog}

La gestión del backlog se realizó mediante GitHub Projects, utilizando un tablero
Kanban asociado al repositorio del proyecto. En este tablero se registraron las
tareas planificadas, incidencias, mejoras y ajustes detectados durante los
sprints. La organización visual por columnas permitió conocer el estado de cada
elemento de trabajo y mantener una visión compartida del avance del proyecto.

Para que las tareas fueran homogéneas y trazables, se definió una plantilla de
GitHub Issues para nuevas tareas dentro del repositorio. Esta plantilla incluía
campos como resumen, módulo afectado, tipo de tarea, prioridad, sprint,
estimación, rama sugerida, dependencias y Definition of Done. Además, se siguió
una identificación organizada de las tareas, reflejada también en los nombres de
ramas del repositorio mediante patrones como \path{feature/tXX-descripcion} o
\path{fix/bugXX-descripcion}. Esta convención facilitó relacionar cada issue con
su rama, su Pull Request y el incremento funcional correspondiente.

\begin{figure}[H]
\centering
\fbox{%
\parbox{0.86\textwidth}{%
\centering
\vspace{1.2cm}
Captura pendiente: tablero Kanban de GitHub Projects con columnas de seguimiento
y tareas del proyecto.
\vspace{1.2cm}
}}
\caption{Tablero Kanban utilizado para la gestión del backlog.}
\label{fig:github_project_kanban}
\end{figure}

\begin{figure}[H]
\centering
\fbox{%
\parbox{0.86\textwidth}{%
\centering
\vspace{1.2cm}
Captura pendiente: plantilla de issue para tareas
\texttt{.github/ISSUE\_TEMPLATE/task.yml}.
\vspace{1.2cm}
}}
\caption{Plantilla de tarea utilizada para homogeneizar el backlog.}
\label{fig:template_tareas_github}
\end{figure}

\subsection{Gestión documental}

La gestión documental combinó varias herramientas según la fase del proyecto. Al
inicio, la memoria se redactó en Overleaf, lo que facilitaba la edición colaborativa
en LaTeX y la revisión rápida del documento. Sin embargo, a medida que la memoria
creció en tamaño, número de capítulos, tablas y anexos, aparecieron limitaciones de
compilación. Por este motivo, la documentación principal se trasladó a un
repositorio de GitHub, manteniendo el control de versiones de los ficheros
\texttt{.tex}, figuras, anexos y PDF generado.

Además, se creó una carpeta compartida en OneDrive para conservar documentos
iniciales de planificación y trabajo, elaborados principalmente con Microsoft Word.
Esta carpeta sirvió como espacio de apoyo para materiales previos, borradores,
actas o documentos que todavía no formaban parte de la memoria final en LaTeX.

\begin{figure}[H]
\centering
\fbox{%
\parbox{0.86\textwidth}{%
\centering
\vspace{1.2cm}
Captura pendiente: estructura documental en Overleaf, repositorio de memoria en
GitHub o carpeta compartida en OneDrive.
\vspace{1.2cm}
}}
\caption{Herramientas utilizadas para la gestión documental del proyecto.}
\label{fig:gestion_documental}
\end{figure}

\subsection{Gestión del código fuente}

El código fuente de OmniLitter se gestionó mediante un repositorio centralizado en
GitHub dentro de la organización \texttt{Omnilitter-TFG}. El repositorio principal,
\texttt{OmniLitter}, agrupó las distintas partes del sistema: backend, portal web,
aplicación móvil, paquete compartido, configuración de despliegue y documentación
técnica auxiliar. Esta estructura permitió trabajar con una visión de monorepo y
mantener en un mismo espacio los contratos compartidos entre API, web y móvil.

Para el control de versiones se utilizó Git. El flujo de trabajo se apoyó en ramas
por funcionalidad o corrección, con prefijos como \texttt{feature/} y \texttt{fix/}.
Las ramas se integraban mediante Pull Requests hacia \texttt{develop} y, una vez
acumulado un incremento estable, hacia \texttt{main}. En el historial del repositorio
se observa este patrón mediante múltiples merges de Pull Request desde ramas como
\path{feature/t21-sync-batch}, \path{feature/t34-entorno-sesion-y-rangos-tamano}
o \path{fix/bug07-graficas-dashboard-movil}.

También se aplicó una convención de mensajes próxima a Conventional Commits, usando
prefijos como \texttt{feat:}, \texttt{fix:}, \texttt{test:} o \texttt{docs:}. Esta
práctica facilitó interpretar el historial, distinguir nuevas funcionalidades de
correcciones y relacionar los cambios con tareas del backlog.

\begin{figure}[H]
\centering
\fbox{%
\parbox{0.86\textwidth}{%
\centering
\vspace{1.2cm}
Captura pendiente: repositorio GitHub de la organización
\texttt{Omnilitter-TFG/OmniLitter}, ramas o historial de Pull Requests.
\vspace{1.2cm}
}}
\caption{Repositorio centralizado utilizado para la gestión del código fuente.}
\label{fig:repositorio_codigo_fuente}
\end{figure}

\subsection{Gestión de versiones y releases}

La evolución del producto se documentó mediante versiones asociadas a incrementos
del desarrollo. Se generaron releases intermedias por sprint, utilizando etiquetas
como \texttt{v0.1.0}, \texttt{v0.2.0} y \texttt{v0.3.0}, y una versión final
\texttt{v.1.0.0}, correspondiente al cierre funcional del proyecto. Estas versiones
permitieron marcar puntos estables del repositorio y relacionar cada entrega con
un conjunto de funcionalidades ya integradas.

Las releases se gestionaron con herramientas de GitHub. El repositorio incluye un
workflow de GitHub Actions que crea una release al publicar una etiqueta que
coincida con el patrón \texttt{v*}, generando notas básicas a partir de commits y
Pull Requests. Además, se incorporaron otros flujos de integración y despliegue,
como validación continua, construcción de imágenes Docker y generación manual de
una build Android de previsualización con Expo/EAS.

\begin{table}[H]
\centering
\caption{Versiones principales publicadas durante el desarrollo}
\label{tab:versiones_releases}
\small
\renewcommand{\arraystretch}{1.18}
\begin{tabularx}{\textwidth}{p{2.2cm}p{3.2cm}X}
\toprule
\textbf{Versión} & \textbf{Momento} & \textbf{Finalidad} \\
\midrule
\texttt{v0.1.0} & Release intermedia & Marcar una primera base funcional integrada del sistema. \\
\texttt{v0.2.0} & Release intermedia & Consolidar avances de captura de campo, persistencia inicial y módulos web asociados. \\
\texttt{v0.3.0} & Release intermedia & Recoger funcionalidades avanzadas de sincronización, mapas, análisis y exportación. \\
\texttt{v.1.0.0} & Release final & Señalar la versión final del proyecto tras el cierre funcional, refinamiento y validación. \\
\bottomrule
\end{tabularx}
\end{table}

\begin{figure}[H]
\centering
\fbox{%
\parbox{0.86\textwidth}{%
\centering
\vspace{1.2cm}
Captura pendiente: sección de releases de GitHub con las etiquetas
\texttt{v0.1.0}, \texttt{v0.2.0}, \texttt{v0.3.0} y \texttt{v.1.0.0}.
\vspace{1.2cm}
}}
\caption{Releases publicadas mediante GitHub.}
\label{fig:github_releases}
\end{figure}

\subsection{Gestión de incidencias}

La gestión de incidencias se integró en el mismo sistema de seguimiento utilizado
para el backlog. Los errores detectados durante el desarrollo o durante las pruebas
manuales se registraron en GitHub Issues y se visualizaron en el tablero Kanban de
GitHub Projects. De esta forma, las incidencias podían priorizarse junto con el
resto de tareas, asignarse a un sprint y seguirse hasta su resolución.

Para homogeneizar el reporte de errores se creó una plantilla específica de bug en
el repositorio. Esta plantilla recogía la descripción del problema, pasos para
reproducirlo, comportamiento esperado, comportamiento actual, módulo afectado,
sprint y severidad. También incluía una Definition of Done orientada a confirmar
que el error había sido corregido, probado localmente, revisado y fusionado en
\texttt{develop}. El repositorio conserva ramas de corrección con nombres como
\path{fix/bug01-responsive-portal}, \path{fix/bug02-roles-web} o
\path{fix/bug07-graficas-dashboard-movil}, lo que muestra la relación entre
incidencia, rama y Pull Request.

\begin{figure}[H]
\centering
\fbox{%
\parbox{0.86\textwidth}{%
\centering
\vspace{1.2cm}
Captura pendiente: plantilla de incidencia
\texttt{.github/ISSUE\_TEMPLATE/bug.yml} o ejemplo de bug en el tablero Kanban.
\vspace{1.2cm}
}}
\caption{Plantilla y seguimiento de incidencias en GitHub.}
\label{fig:template_incidencias_github}
\end{figure}

